\documentclass[12pt,letterpaper]{article}



%%
%% Required packages:
%%
\usepackage{etex}
\usepackage{amsmath}
\usepackage{amssymb}
\usepackage{amstext}
\usepackage{amsthm}
\usepackage{fancyhdr,fancybox}
\usepackage{mathrsfs} % need for \mathscr command
\usepackage{multicol}
\usepackage[normalem]{ulem}
%\usepackage{pictex,pictexplus}
% \usepackage{pictexwd}
\usepackage{rotating}
\usepackage{url}
\usepackage{xfrac}
\usepackage{booktabs}
\usepackage{color,soul}
\usepackage{comment}

\usepackage{hyperref}
\hypersetup{
    colorlinks=true,     % false: boxed links; true: colored links
    linkcolor=blue,      % color of internal links
    citecolor=blue,      % color of links to bibliography
    filecolor=blue,      % color of file links
    urlcolor=blue        % color of external links
}


\usepackage{tikz}
\usetikzlibrary{backgrounds,calc}

%%
%% Common formatting commands
%%
\setlength{\parindent}{0in}
\setlength{\textwidth}{7in}
\setlength{\evensidemargin}{-0.25in}
\setlength{\oddsidemargin}{-0.25in}
\setlength{\parskip}{.5\baselineskip}
\setlength{\topmargin}{-0.5in}
\setlength{\textheight}{9in}

%               Problem and Part
\newcounter{problemnumber}
\newcounter{partnumber}
\newcounter{subpartnumber}
\newcommand{\Problem}{%
  \refstepcounter{problemnumber}%
  \setcounter{partnumber}{0}%
  \setcounter{subpartnumber}{0}%
  \item[\makebox{\hfil\textbf{\theproblemnumber.}\hfil}]%
  }
\newcommand{\InProblem}[2][1.6]{%
  \refstepcounter{problemnumber}%
  \setcounter{partnumber}{0}%
  \setcounter{subpartnumber}{0}%
  \textbf{\theproblemnumber.}\ \ \parbox[t]{#1in}{#2}%
}
\newcommand{\InSmallProblem}[1]{\InProblem[1.05]{#1}}
\newcommand{\NoProblem}[1]{%
  \mbox{\hfil\textbf{#1}\hfil}%
  }
\newcommand{\UnProblem}{%
  \refstepcounter{problemnumber}%
  \setcounter{partnumber}{0}%
  \setcounter{subpartnumber}{0}%
  \mbox{\hfil\textbf{\theproblemnumber.}\hfil}%
  }
\newcommand{\InFlexProblem}[1]{%
  \refstepcounter{problemnumber}%
  \setcounter{partnumber}{0}%
  \setcounter{subpartnumber}{0}%
  \textbf{\theproblemnumber.}\ \ #1%
}
%%  Part:
\newcommand{\Part}{%
  \renewcommand{\thepartnumber}{(\alph{partnumber})}%
  \refstepcounter{partnumber}
  \setcounter{subpartnumber}{0}%
  \item[(\alph{partnumber})]%
}
\newcommand{\InPart}[2][1.75]{%
  \renewcommand{\thepartnumber}{(\alph{partnumber})}%
  \refstepcounter{partnumber}%
  \setcounter{subpartnumber}{0}%
  (\alph{partnumber})\ \ \parbox[t]{#1in}{#2}%
}
\newcommand{\UnPart}{%
  \refstepcounter{partnumber}%
  \renewcommand{\thepartnumber}{(\alph{partnumber})}%
  \setcounter{subpartnumber}{0}%
  (\alph{partnumber})%
}
\newcommand{\InLargePart}[1]{\InPart[3.05]{#1}}
\newcommand{\InFullPart}[1]{\InPart[6.2]{#1}}
\newcommand{\InSmallPart}[1]{\InPart[1.05]{#1}}
%% SubPart:
\newcommand{\SubPart}{%
  \refstepcounter{subpartnumber}%
  \item[(\roman{subpartnumber})]%
  }
\newcommand{\InSubPart}[2][2.25]{%
  \stepcounter{subpartnumber}%
  (\roman{subpartnumber})\ \ \parbox[t]{#1in}{#2}%
}
\newcommand{\InSmallSubPart}[1]{\InSubPart[1.5]{#1}}
\newcommand{\InTinySubPart}[1]{%
  \stepcounter{subpartnumber}%
  (\roman{subpartnumber})\ \ \makebox{#1}%
}
\newcommand{\InMySubPart}[2][2.25]{%
  \stepcounter{subpartnumber}%
  \parbox[t]{#1in}{\makebox[0.3in][l]{(\roman{subpartnumber})}\ \ {#2}}%
}
\newcommand{\TrueFalse}{\textbf{T}\ \ \textbf{F}\ \ }
\newcommand{\TFPart}[1]{% 
  \stepcounter{partnumber}%
  \setcounter{subpartnumber}{0}%
  \item[(\alph{partnumber})] \TrueFalse\parbox[t]{5.5in}{#1}}

\newcommand{\Points}[1]{(#1 points) \ }
\newcommand{\Point}[1]{(#1 point) \ }


%% For Answers:
\newcounter{answernumber}
\newcommand{\Answer}{%
  \refstepcounter{answernumber}%
  \setcounter{partnumber}{0}%
  \setcounter{subpartnumber}{0}%
  \item[\makebox{\hfil\textbf{\theanswernumber.}\hfil}]%
  }


% Setting up the headers for the first page
%\chead{} % will default to what is typed in the file.
\fancyhead[L]{\textsc{Math 22a}}
\fancyhead[R]{\textsc{Fall 2024}}
\fancyfoot[R]{
	\vspace*{\fill}\footnotesize\textcopyright{} \emph{Harvard Math 22a}	
}
\renewcommand{\headrulewidth}{0pt}

%\fancyhead[C]{}
%	\chead{}	
%	\lhead{}
%	\rhead{}
%	\cfoot{}


% Putting copyright on the footer of each page
%\fancypagestyle{secondpage}{
%	\fancyhead[C]{TESTing again} % change this to be blank
%		%\chead{}	
%	\fancyhead[L]{bears} % change this to be blank
%	\fancyhead[R]{dogs} % change this to be blank
%	\fancyfoot[R]{ %keeps the footer the same
%		\vspace*{\fill}\footnotesize\textcopyright{} \emph{Harvard Math 22a}	
%	}
%}
%\renewcommand{\headrulewidth}{0pt}
%\pagestyle{fancy}
\pagestyle{empty}


%\lhead{}
%\rhead{}
%\rfoot{\vspace*{\fill}\footnotesize\textcopyright{} \emph{Harvard Math 22a}}
%\renewcommand{\headrulewidth}{0pt}
%\pagestyle{fancy}




%%
%% Common shortcut commands:
%%
\newcommand{\ds}{\displaystyle}
\newcommand{\sgn}{\operatorname{sgn}}
\renewcommand{\Pr}{\operatorname{P}}
\newcommand{\CPr}[3][{}]{\Pr_{\text{#1}}\left(#2\,|\,#3\right)}

\newcommand{\C}{\mathbb{C}}
\newcommand{\R}{\mathbb{R}}
\newcommand{\Z}{\mathbb{Z}}
\newcommand{\N}{\mathbb{N}}
\newcommand{\Q}{\mathbb{Q}}
\newcommand{\D}{\mathbb{D}}

\newcommand{\rref}{\operatorname{rref}}
\newcommand{\im}{\operatorname{im}}
\newcommand{\rank}{\operatorname{rank}}
\newcommand{\diag}{\operatorname{diag}}
\newcommand{\Span}{\operatorname{span}}
%\renewcommand{\vec}[1]{\mathbf{#1}}
\newcommand{\charpoly}{\mathscr{P}}
\newcommand{\nicefrac}[2]{\sfrac{#1}{#2}} % using xfrac package
\newcommand{\topology}{\mathcal{T}}
\newcommand{\Inn}{\operatorname{Inn}}

\newcommand{\blankbox}[2]{\fbox{\rule{#1}{0in}\rule{0in}{#2}}}

\newenvironment{myindentpar}[1]%
 {\begin{list}{}%
         {\setlength{\leftmargin}{#1}
          \setlength{\rightmargin}{\leftmargin}}%
         \item[]%
 }
 {\end{list}}
\newtheorem{theorem}{Theorem}
\newtheorem*{NStheorem}{Theorem}
\newtheorem{cor}[theorem]{Corollary}
\newtheorem*{claim}{Claim}
\newtheorem{defn}{Definition}

\newenvironment{amatrix}[1]{%
  \left(\begin{array}{@{}*{#1}{c}|c@{}}
}{%
  \end{array}\right)
}

%--------grstep
% For denoting a Gauss' reduction step.
% Use as: \grstep{\rho_1+\rho_3} or \grstep[2\rho_5 \\ 3\rho_6]{\rho_1+\rho_3}
\newcommand{\grstep}[2][\relax]{%
   \ensuremath{\mathrel{
       {\mathop{\longrightarrow}\limits^{#2\mathstrut}_{
                                     \begin{subarray}{l} #1 \end{subarray}}}}}}
\newcommand{\swap}{\leftrightarrow}


\usetikzlibrary{arrows}
\usepackage{wrapfig}
\usepackage{amsfonts,amsmath}

\makeatletter
\renewcommand*\env@matrix[1][*\c@MaxMatrixCols c]{%
	\hskip -\arraycolsep
	\let\@ifnextchar\new@ifnextchar
	\array{#1}}
\makeatother

\def\env@matrix{\hskip -\arraycolsep
	\let\@ifnextchar\new@ifnextchar
	\array{*\c@MaxMatrixCols c}}

\newcommand{\norm}[1]{\left\lVert#1\right\rVert}

\chead{\Large\textbf{Least Squares, $\vec\S$6.5, 6.6}}% 

\begin{document}
\thispagestyle{fancy}

Recall from previously:

\begin{itemize}
\item An {\bf orthogonal matrix} $U$ is a square matrix so that $U^{-1}=U^{T}$.
\item  {\bf Orthogonal Decomposition Theorem:} Let $W$ be a subspace of $\mathbb{R}^n$. Then each $y \in \mathbb{R}^n$ can be uniquely written as $$y=\hat{y} + z$$ where $\hat{y} \in W$ and $z \in W^{\perp}$. 

Moreover, if $\{u_1, \dots, u_p\}$ is any orthogonal basis for $W$, then $$\hat{y}= \left(\frac{y \cdot u_1}{u_1 \cdot u_1} \right)u_1+\dots+  \left(\frac{y \cdot u_p}{u_p \cdot u_p} \right)u_p,$$ and $z=y-\hat{y}$.
\item The vector $\hat{y}$ is the {\bf orthogonal projection} of $y$ onto the subspace $W$. Written as $$\hat{y}=\text{proj}_{W}{y}.$$
\item {\bf Best Approximation Theorem:} Let $W$ be a subspace of $\mathbb{R}^n$. Let $\hat{y}$ be the orthogonal projection of $y$ onto $W$, then $$\norm{y-\hat{y}} < \norm{y-v}$$ for all $v\in W \setminus \{\hat{y}\}$.
\item {\bf Theorem 6.10:} If $\{ u_1, \dots, u_p\}$ is an orthonormal basis for $W$, then $$\text{proj}_{W}{y} =(y\cdot u_1) u_1 +\dots +(y \cdot u_p) u_p.$$ If $U=[u_1 \; \dots \; u_p$, then $$\text{proj}_{W}{y} =U U^{T} y.$$
\item {\bf Gram-Schmidt Orthogonalization.} Given a basis $\{x_1, \dots, x_p\}$ for a non-zero subspace $W$ of $\mathbb{R}^n$, define 
\begin{align*}
v_1 & = x_1\\
v_2 & = x_2 - \left( \frac{x_2 \cdot v_1}{v_1 \cdot v_1} \right) v_1\\
v_3 & = x_3 -\left( \frac{x_3 \cdot v_1}{v_1 \cdot v_1} \right) v_1-\left( \frac{x_3 \cdot v_2}{v_2 \cdot v_2} \right) v_2\\
&\vdots\\
v_p & = x_p -\left( \frac{x_p \cdot v_1}{v_1 \cdot v_1} \right) v_1-\dots -\left( \frac{x_p \cdot v_{p-1}}{v_{p-1} \cdot v_{p-1}} \right) v_{p-1}.
\end{align*}
Then $\{v_1, \dots, v_p\}$ is an orthogonal basis for $W$. In addition $\text{Span}\{v_1,\dots, v_k\}=\text{Span}\{ x_1, \dots, x_k\}$ for all $1 \leq k \leq p$.
\item {\bf QR Factorization.} If $A$ is an $m\times n$ matrix with linearly independent columns, then $A=QR$ where $Q$ is an $m\times n$ matrix whose columns form an orthonormal basis for the column space of $A$ and $R$ is an upper triangular $n \times n$ matrix with positive entries along the diagonal. 
\end{itemize}

\newpage


{\bf Goal:} Solve the matrix $Ax=b$.

What if the matrix equation is inconsistent?

\begin{defn} A {\bf least squares solution} to $Ax=b$ is a vector $\hat{x} \in \mathbb{R}^n$ such that $$\norm{b-A\hat{x}} \leq \norm{b-Ax}$$ for all $x \in \mathbb{R}^n$.
\end{defn}

How can we find least squares solutions?

{\bf Theorem 6.13:} The set of least squares solutions to $Ax=b$ is the same as the set of solutions to the normal squares $A^{T}Ax=A^{T}b$.

\begin{enumerate}



\Problem Prove the theorem.

\vfill

\newpage

\begin{defn} The {\bf least squares error} is given by the distance from $b$ to $A\hat{x}$.
\end{defn}

\Problem Let $A=\begin{bmatrix} 4 & 0\\ 0 & 2\\ 1 & 1\\ \end{bmatrix}$ and $b=\begin{bmatrix} 2\\ 0\\ 11\\ \end{bmatrix}$. Find a least squares solutions to $Ax=b$.

\vfill

\Problem Let $A=\begin{bmatrix} 1 & -6\\ 1 & -2\\ 1 & 1\\1 & 7\\ \end{bmatrix}$ and $b=\begin{bmatrix} -1\\ 2\\ 1\\ 6\\ \end{bmatrix}$. Find a least squares solutions to $Ax=b$.

\vfill

\newpage

\Problem Find the best fit line through the points $(-6,-1), (-2,2), (1,1), (7,6).$

\vfill

\Problem   Here are the number of daylight hours in Cambridge on a few days of
  2019, from the Astronomical Applications Department of the
  U.S. Naval Observatory.
  % https://aa.usno.navy.mil/cgi-bin/aa_rstablew.pl?ID=AA&year=2019&task=0&state=MA&place=cambridge

  \begin{center}
    \begin{tabular}{l l l}
      Day & Which day of the year & number of daylight hours \\
      \hline
      May 17 & 137th & 14 hours and 40 minutes \\
      June 1 & 152nd & 15 hours and 5 minutes \\
      July 1 & 182nd & 15 hours and 14 minutes \\
      August 8 & 220th & 14 hours and 12 minutes
    \end{tabular}
  \end{center}

  Suppose you would like to fit a function $f$ that gives the
  number of hours of daylight on a particular day, so that you can
  predict the number of daylight hours on the shortest day of the
  year.  Describe how you would go about this; what form would your
  function $f$ have? 
  
  \vfill

\newpage

\end{enumerate}

\begin{comment}

\newpage
\chead{\Large\textbf{Least Squares -- Answers / Solutions}}%
\lhead{}
\rhead{}
\thispagestyle{fancy}

\begin{enumerate}
\Answer 
\begin{proof}
We prove the two sets are equal by showing that each is a subset of the other.

First, let $\hat{x}$ be a least squares solutions to $Ax=b$. Let $\hat{b} =\text{proj}_{\text{Col(A)}}{b}$. Thus $A\hat{x} =\hat{b}$. By the Orthogonal Decomposition Theorem, $b -\hat{b}$ is orthogonal to $\text{Col}(A)$, and hence $b-\hat{b}$ is orthogonal to the columns $a_j$ of $A$. Thus we have $a_j^{T}(b-A\hat{x})=0$ for $1\leq j \leq n$. Hence $A^{T}(b-A\hat{x})=0$, and rearranging we get $A^T b =A^{T} A \hat{x}$. Thus $\hat{x}$ is a solution to the normal equations.

Second, suppose $\hat{x}$ is a solution to the normal equations. Rearranging the normal equations, we get $$A^{T}(b-A\hat{x})=0.$$ Thus $b-A\hat{x}$ is orthogonal to the the column space of $A$. Since $b=A\hat{x} +(b-A\hat{x})$ is an orthogonal decomposition of $b$ and the orthogonal decomposition is unique, we know $A\hat{x}=\text{proj}_{\text{Col}(A)}{b}$. By the Best Approximation Theorem, $\hat{x}$ is a least squares solution.
\end{proof}
\Answer We use the normal equations to find the least squares solutions.
Note $$A^{T}A=\begin{bmatrix} 17 & 1\\ 1 & 5\\ \end{bmatrix}\text{ and } A^{T}b=\begin{bmatrix} 19\\ 1\\ \end{bmatrix}.$$ Thus $$\hat{x} = \begin{bmatrix} 17 & 1\\ 1 & 5\\ \end{bmatrix}^{-1} \begin{bmatrix} 19\\11\\ \end{bmatrix}=\begin{bmatrix} 1\\ 2\\ \end{bmatrix}.$$
\Answer Note the columns of $A$ are orthogonal. Now $$\hat{b}=\text{proj}_{\text{Col}(A)}{b} =\frac{b \cdot a_1}{a_1 \cdot a_1} a_1 +\frac{b \cdot a_2}{a_2\cdot a_2} a_2=\frac{8}{4}a_1+\frac{45}{90}a_2=\begin{bmatrix} -1\\ 1\\ 5/2\\ 11/2\\ \end{bmatrix}.$$ Now we need to solve $Ax =\hat{b}$, and hence we get $\hat{x}=\begin{bmatrix} 2\\ 1/2\\ \end{bmatrix}$.
\Answer We fit the data to the line $y=\beta_0 +\beta_1 x$, and hence we get the following four equations.
\begin{align*}
-1  =& \beta_0 +\beta_1 (-6)\\
2  =& \beta_0 +\beta_1 (-2)\\
1  =& \beta_0 +\beta_1 (1)\\
6  =& \beta_0 +\beta_1 (7)\\
\end{align*}
Thus we can use the matrix $A$ and the vector $b$ from the last problem. The least squares solution is then $\begin{bmatrix} 2\\ 1/2\\ \end{bmatrix}$. Thus the best fit line is $y=2+\frac{x}{2}$.
\Answer We can try to fit the data to the function $$f(t)=c_0+c_1 \cos \left( \frac{2\pi}{365} t\right)+c_2 \sin \left( \frac{2\pi}{365} t\right).$$
\end{enumerate}  

\end{comment}
\end{document}


%%% Local ables: 
%%% mode: latex
%%% TeX-master: t
%%% End: 
